\DocSection{Introduction}
\label{sec:introduction}

Planning in a stochastic environment requires two conceptually separate
operations. First, a policy is evaluated by averaging its return over the
environment's transition randomness. Second, a decision rule selects among
policies. Conflating these operations changes the objective: exponentiating a
single realized return before averaging is generally not equivalent to
exponentiating the expected return.

This manuscript develops a policy-inference formulation in which the target is
defined before an approximation algorithm is chosen. The target assigns mass
to complete deterministic policies according to expected return. Its modes are
therefore expected-return maximizers, while a finite inverse temperature
retains a distribution over alternatives. The corrected inference method uses
rollouts only through estimators that are valid for that target.

The study has four goals:
\begin{enumerate}
  \item specify the expected-return policy distribution and its assumptions;
  \item construct exact diagnostic oracles that detect target substitution;
  \item derive direct and particle estimators with a consistent execution rule; and
  \item compare target fidelity, control performance, and computational cost
        under matched experimental conditions.
\end{enumerate}

Agreement with exact enumeration establishes target correctness before any
large-scale performance comparison is attempted.
